% ----- Consignes exo 4 ----- %
\begin{td-exo}[Autour de \textsc{Coupe maximum}]\, % 4 
	\begin{problem}{\textsc{Coupe maximum} (CUT)}
    \Input{Soit \(G = (V, E)\) un graphe non orienté et \(k\in\bb N\).}
    \Question{Existe-t-il une partition des sommets en deux sous-ensembles \(V_1\) et \(V_2\) tels que le nombre d'arêtes entre \(V_1\) et \(V_2\) soit \(k\)?}
\end{problem}

	Réduire \textsc{Non Egal 3-SAT} à \textsc{Coupe maximum}. Conclure.
	
	\textit{Remarque}: vous verrez en master que coupe min est polynomiale meme si les aretes ont des poids.
\end{td-exo}

% ----- Solutions exo 4 ----- %
\iftoggle{showsolutions}{ 
	\begin{td-sol}[]\ % 4
		On supposera que dans les clauses de taille \(3\) on a pas à la fois \(x\) et \(\ol{x}\), sinon la clause est toujours satisfaite et on peut l'ignorer.
		On suppose de plus que toute paire de clauses a au plus un littéral en commun. On procède alors à la construction suivante:
		\begin{constr}
			On crée un graphe \(G\) qui, à chaque littéral, crée un sommet \(v\) et un sommet \(\ol{v}\).
			% voir construction

			On obtient alors que le nombre d'arêtes sur ce graphe est égal à 3 fois le nombre de clauses de \(\varphi\) plus une par variable, soit \(3m + n\).
			De plus, on a
			\begin{equation*}
				\sum_{x_i \in L} n_{x_i} = 3m
			\end{equation*}
			Alors, on cherche une coupe telle que
			\begin{equation*}
				C(S, \ol{S}) \geq 2m + n
			\end{equation*}

			Soit \(t\) une solution de \textsc{Non Egal 3-SAT}. Supposons que \(n\) variables sont mises à vrai, et donc \(n\) variables sont mises à faux. Alors, on a \(2m\) arêtes de clauses coupées et \(n\) arêtes de variables coupées, soit un total de \(2m + n\) arêtes coupées.
			Plus précisément, avec les \(2m\) arêtes restantes, on ne peut couper que des triangles. Ces triangles, ayant au moins 1 sommet dans chaque partie de la coupe, sont forcément coupés 2 fois. Il y en a \(m\) au total, soit \(2m\) arêtes coupées.
		\end{constr}
	\end{td-sol}
}{}
